% !PRINTED 470
\weberpage{470}

%<<CONT
\[
\varphi(\varrho) = \varphi(\varrho_a)
\]
then we say that \gesperrt{$\omega$ permits or admits the substitution $(\varrho, \varrho_a)$ or $\sigma_a$}.

The quantities $\omega$ which, apart from the identical substitution, admit no substitution $\sigma_a$, are the primitive elements of the field $\Omega(\varrho)$.

\gesperrt{4. A quantity which admits all the substitutions $\sigma_a$ is necessarily an element of $\Omega$.}

For such a quantity is identical with all its conjugate quantities, and hence by \S.\,144 satisfies an equation of the first degree in $\Omega$.

\begin{center}
\S.\ 147.
\end{center}

\begin{center}
\gesperrt{Composition of the Substitutions.}
\end{center}

If in any function of $\varrho$ we replace the root $\varrho$ first by $\varrho_a$ and then $\varrho_a$ by $\varrho_b$, the result is the same as if we had from the outset interchanged $\varrho$ with $\varrho_b$. If therefore we set
\weq{1}{\sigma = (\varrho, \varrho_a), \quad \sigma' = (\varrho_a, \varrho_b), \quad \sigma'' = (\varrho, \varrho_b),}
then it is indifferent for the result whether in all the numbers of the field $\Omega(\varrho)$ we first carry out the substitution $\sigma$ and then the substitution $\sigma'$, or whether we make the substitution $\sigma''$ in a single step.

\gesperrt{1. We therefore call $\sigma''$ compounded or composed of $\sigma$ and $\sigma'$, and denote this relation by the symbolic equation}
\weq{2}{\sigma \sigma' = \sigma''.}
Since by Theorem 3.\ of the preceding section we may, in the designation of a substitution $(\varrho_h, \varrho_k)$, choose the first or the second element arbitrarily, $\varrho_a, \varrho_b$ can be so selected that $\sigma, \sigma'$ are any two arbitrarily given ones among the $\mu$ substitutions of the field $\Omega(\varrho)$; but $\sigma''$ is thereby completely determined. But likewise, if $\sigma$ and $\sigma''$ are given, $\sigma'$ is uniquely determined. For if we choose $\varrho$ arbitrarily, then $\varrho_a$ is completely determined by $\sigma$, and $\varrho_b$ by $\sigma''$, whereby $\sigma' = (\varrho_a, \varrho_b)$ is also given. Finally, if $\sigma'$ and $\sigma''$ are given, then $\sigma$ is uniquely determined, since first $\varrho_b$ is determined by $\sigma''$ and then $\varrho_a$ by $\sigma'$. We therefore have:
%CONT>>
